% Mapa declarativo de retitulado aprobado. No altera los cuerpos científicos.
% Cada entrada procede de una coincidencia estática uno-a-uno y conserva el nivel.
\makeatletter
\DeclareUnicodeCharacter{00B2}{\ensuremath{^{2}}}
\DeclareUnicodeCharacter{00B3}{\ensuremath{^{3}}}
\DeclareUnicodeCharacter{00B7}{\ensuremath{\cdot}}
\DeclareUnicodeCharacter{00D7}{\ensuremath{\times}}
\DeclareUnicodeCharacter{0393}{\ensuremath{\Gamma}}
\DeclareUnicodeCharacter{03B1}{\ensuremath{\alpha}}
\DeclareUnicodeCharacter{03B6}{\ensuremath{\zeta}}
\DeclareUnicodeCharacter{03BE}{\ensuremath{\xi}}
\DeclareUnicodeCharacter{2013}{--}
\DeclareUnicodeCharacter{2074}{\ensuremath{^{4}}}
\DeclareUnicodeCharacter{2076}{\ensuremath{^{6}}}
\DeclareUnicodeCharacter{207B}{\ensuremath{^{-}}}
\DeclareUnicodeCharacter{2080}{\ensuremath{_{0}}}
\DeclareUnicodeCharacter{2081}{\ensuremath{_{1}}}
\DeclareUnicodeCharacter{2082}{\ensuremath{_{2}}}
\DeclareUnicodeCharacter{2083}{\ensuremath{_{3}}}
\DeclareUnicodeCharacter{2085}{\ensuremath{_{5}}}
\DeclareUnicodeCharacter{2087}{\ensuremath{_{7}}}
\DeclareUnicodeCharacter{2088}{\ensuremath{_{8}}}
\DeclareUnicodeCharacter{2089}{\ensuremath{_{9}}}
\DeclareUnicodeCharacter{210F}{\ensuremath{\hbar}}
\DeclareUnicodeCharacter{2192}{\ensuremath{\to}}
\DeclareUnicodeCharacter{21A6}{\ensuremath{\mapsto}}
\DeclareUnicodeCharacter{2212}{\ensuremath{-}}
\DeclareUnicodeCharacter{2218}{\ensuremath{\circ}}
\DeclareUnicodeCharacter{221A}{\ensuremath{\surd}}
\DeclareUnicodeCharacter{2260}{\ensuremath{\ne}}
\DeclareUnicodeCharacter{2261}{\ensuremath{\equiv}}
\DeclareUnicodeCharacter{27E8}{\ensuremath{\langle}}
\DeclareUnicodeCharacter{27E9}{\ensuremath{\rangle}}
\newcommand{\HMTTitleMapEntry}[2]{%
  \edef\HMTTitleProbe{\detokenize{#1}}%
  \ifx\HMTTitleKey\HMTTitleProbe\def\HMTResolvedTitle{#2}\fi}
\newcommand{\HMTResolvePublicTitle}[1]{%
  \def\HMTResolvedTitle{#1}%
  \edef\HMTTitleKey{\detokenize{#1}}%
  % El prefacio fue aprobado como texto autoral literal; su subtítulo no se
  % retitula ni se normaliza.
  \HMTTitleMapEntry{Interaction complexe entre deux machines simples}{Interaction complexe entre deux machines simples}% ADP102-0001-LITERAL
  \HMTTitleMapEntry{Introduction générale}{Architecture causale du problème du continu et portée mathématico-physique de HMT–MD}% ADP102-0002
  \HMTTitleMapEntry{Objet initial et ordre causal : échelle finie, horizon d'observation et publication archimédienne}{Horizon projectif : dépendance entre échelle, bord et mesure}% ADP102-0003
  \HMTTitleMapEntry{Intervalle orienté, marques, cellules et arête de clôture}{Construction de la droite finie par l'ordre, l'échelle et la conservation de l'unité}% ADP102-0004
  \HMTTitleMapEntry{Raffinement, survie et seuils de prolongement compatible}{Relation de limite entre infinitésimal, prolongement infini et seuil projectif}% ADP102-0005
  \HMTTitleMapEntry{Évaluation archimédienne des sections compatibles et reconstruction de \(\mathbb R\)}{Limite compatible d'horizons finis et réalisation de la droite réelle}% ADP102-0006
  \HMTTitleMapEntry{Noyau formel de l'holographie modulaire triadique}{Construction et architecture opératorielle du noyau formel APP–TRIT–TPK : support arithmétique, orientation ternaire et dynamique généalogique}% ADP102-0007
  \HMTTitleMapEntry{Relèvements et cocycle de retenue}{Extensions résidu–quotient et cocycle nonadique de retenue}% ADP102-0008
  \HMTTitleMapEntry{Bloc central de \(2\) positions}{Bloc central bivarié d'APP et décomposition en 2 modes spectraux}% ADP102-0009
  \HMTTitleMapEntry{Courbure mixte des \(2\) lois}{Courbure discrète de l'interaction entre les feuillets additif et multiplicatif}% ADP102-0010
  \HMTTitleMapEntry{Correspondance entre \(2\) cartes d’une même classe résiduelle}{Conjugaison entre les cartes positive et résiduelle d'une même classe modulo 9}% ADP102-0011
  \HMTTitleMapEntry{Noyau intérieur et complément gnomonique}{Décomposition de l'orthogramme APP en noyau intérieur et complément gnomonique}% ADP102-0012
  \HMTTitleMapEntry{Segments marqués, interstices et cellule nonadique}{Construction de la cellule nonadique à partir de segments marqués et d'interstices}% ADP102-0013
  \HMTTitleMapEntry{Carte positionnelle finie et projection résiduelle décimale}{Projection résiduelle décimale de la carte positionnelle finie d'APP}% ADP102-0014
  \HMTTitleMapEntry{Extraction de la cellule de \(9\) phases}{Construction de la cellule nonadique à 9 phases par projection résiduelle}% ADP102-0015
  \HMTTitleMapEntry{Position, hauteur et division euclidienne}{Coordonnées de position et de hauteur induites par la division euclidienne}% ADP102-0016
  \HMTTitleMapEntry{Domaine inaugural de la cellule nonadique et exclusion causale de l’APP, du TRIT, du TPK et des constantes analytiques}{Domaine arithmétique autonome de la cellule nonadique et précédence causale par rapport à APP–TRIT–TPK et aux constantes analytiques}% ADP102-0017
  \HMTTitleMapEntry{Loi du défaut à module fixe et longueur variable}{Invariance du défaut cyclotomique à module fixé sous variation de longueur}% ADP102-0018
  \HMTTitleMapEntry{Opérateur cyclotomique primaire}{Opérateur cyclotomique d'APP et action sur les 12 fibres de type A₂}% ADP102-0019
  \HMTTitleMapEntry{Localisation de l’agrégat \texorpdfstring{\(54\)}{54}}{Dérivation du défaut somme–produit 459−405=54 dans le bloc central d'APP}% ADP102-0020
  \HMTTitleMapEntry{Résonance de l’opérateur d’incidence en rang \(6\)}{Rang 6 de l'opérateur d'incidence et multiplicité de son mode résonant}% ADP102-0021
  \HMTTitleMapEntry{TRIT comme unité minimale d’information topologique : orientation ternaire, mémoire de retenue et ordre temporel}{TRIT comme unité ternaire orientée : géométries internes, ordre temporel et conditions dynamiques de cristal temporel}% ADP102-0022
  \HMTTitleMapEntry{Composition et cocycle de retenue}{Cocycle de retenue de la composition ternaire orientée}% ADP102-0023
  \HMTTitleMapEntry{Famille universelle d’algèbres quadratiques}{Classification unifiée des 3 algèbres quadratiques du TRIT}% ADP102-0024
  \HMTTitleMapEntry{Topological Phase Kernel : base observable et état enrichi}{TPK : base observable, état enrichi et séparation typée de ses projections}% ADP102-0025
  \HMTTitleMapEntry{La base double des chemins APP}{Catégorie bivariée de chemins APP comme base de l'état enrichi TPK}% ADP102-0026
  \HMTTitleMapEntry{Phase nonadique, secteur dodécaphasique et calendrier}{Extension non scindée C₁₂→C₁₀₈→C₉ et registre dodécaphasique de la phase nonadique}% ADP102-0027
  \HMTTitleMapEntry{Cylindre et frontière}{Compatibilité cylindrique et signature de bord de l'état enrichi}% ADP102-0028
  \HMTTitleMapEntry{Transport des fibres}{Opérateur cardinal \(T_d\) de transport entre fibres enrichies avec conservation de la voie et de la mémoire}% ADP102-0029
  \HMTTitleMapEntry{Relations d’extension}{Relation typée \(\operatorname{Ext}_r\) de prolongement entre fibres enrichies}% ADP102-0030
  \HMTTitleMapEntry{Projections et pertes typées}{Quotients observables de l'état enrichi et champs généalogiques conservés}% ADP102-0031
  \HMTTitleMapEntry{Arbre opératoriel du conteneur}{Insertion des 6 familles fibrées de l'état enrichi dans la hiérarchie opératorielle du TPK}% ADP102-0032
  \HMTTitleMapEntry{Algèbre des opérateurs du Topological Phase Kernel et dynamique causale d’actualisation de l’état enrichi}{Hiérarchie typée des 34 constructions canoniques hétérogènes et dynamique causale d'actualisation du TPK}% ADP102-0033
  \HMTTitleMapEntry{Classification de l’état enrichi en \(8\) classes structurelles invariantes}{Hiérarchie typée L₀–L₇ des 34 constructions canoniques hétérogènes du TPK}% ADP102-0034
  \HMTTitleMapEntry{Décomposition fonctionnelle en \(14\) regroupements compatibles avec les opérateurs APP–TRIT–TPK}{Regroupement fonctionnel des 34 constructions canoniques en 14 familles de composition APP–TRIT–TPK}% ADP102-0035
  \HMTTitleMapEntry{Domaine et typage des \(34\) entrées de l’opérateur d’actualisation}{Classes mathématiques, domaines, codomaines et dépendances causales des 34 constructions canoniques}% ADP102-0036
  \HMTTitleMapEntry{Composition causale de l’actualisation de l’état enrichi}{Actualisation de l'état enrichi au moyen de \(U_t=\operatorname{Upd}_t\circ\operatorname{Tra}_t\circ\operatorname{Sel}_t\)}% ADP102-0037
  \HMTTitleMapEntry{Séparation typée entre opérateurs de transport, de mémoire, de calendrier, d’extraction et de retour}{Séparation typée entre opérateurs, états, projections et publications du TPK}% ADP102-0038
  \HMTTitleMapEntry{Décomposition opératorielle des canaux d’incidence \texorpdfstring{\(90/120\)}{90/120} : sélecteur temporel, générateur nilpotent et composante angulaire}{Incidence tritique 90/120 et séparation typée des opérateurs temporel, nilpotent et angulaire}% ADP102-0039
  \HMTTitleMapEntry{Distribution des indices \(90/120\) entre les \(6\) objets du système d’incidence}{Emplacements typés des indices 90/120 dans 6 objets opératoriels non équivalents}% ADP102-0040
  \HMTTitleMapEntry{Défaut normalisé de l’incidence}{Défaut normalisé 1−90/120=1/4 de l'incidence tritique}% ADP102-0041
  \HMTTitleMapEntry{Invariants de typage}{Invariants de domaine, d'orientation et de mémoire des réalisations 90/120}% ADP102-0042
  \HMTTitleMapEntry{État local et rétroaction de feuillet}{Opérateur local de transition orientée et actualisation du feuillet APP}% ADP102-0043
  \HMTTitleMapEntry{Cylindres des \(2\) lecteurs}{Systèmes cylindriques des 2 lecteurs arithmétiques et compatibilité de troncature}% ADP102-0044
  \HMTTitleMapEntry{Décomposition du Topological Phase Kernel en opérateurs d’avancée, de rotation, de retenue et de mémoire}{Factorisation de la transition TPK sous la forme \(U_t=\operatorname{Upd}_t\circ\operatorname{Tra}_t\circ\operatorname{Sel}_t\)}% ADP102-0045
  \HMTTitleMapEntry{Support ordonné et coordonnée cyclique}{Coordonnée cyclique sur le support temporel orienté à 12 secteurs}% ADP102-0046
  \HMTTitleMapEntry{Extension cyclique non scindée}{Extension non scindée 0→C₁₂→C₁₀₈→C₉→0 et cocycle de retenue}% ADP102-0047
  \HMTTitleMapEntry{Extracteur signé d’événements}{Extraction de l'état dodécaphasique signé à partir du registre temporel enrichi}% ADP102-0048
  \HMTTitleMapEntry{\(2\) mots dodécaphasiques de provenances différentes}{Distinction généalogique entre 2 mots dodécaphasiques de même période}% ADP102-0049
  \HMTTitleMapEntry{Raffinement minimal par microinstants}{Subdivision temporelle minimale compatible avec la semi-conjugaison}% ADP102-0050
  \HMTTitleMapEntry{\(3\) opérations quaternaires non interchangeables}{Commutateurs de 3 opérations quaternaires et séparation de leurs actions}% ADP102-0051
  \HMTTitleMapEntry{Lecteur tritique de l’état orienté : ordre temporel et transport de retenue}{Lecteur tritique orienté avec transport de quotient et de retenue}% ADP102-0052
  \HMTTitleMapEntry{Quotient minimal de mémoire prédictive}{Quotient de mémoire d'ordre 1 et suffisance prédictive de la projection modale}% ADP102-0053
  \HMTTitleMapEntry{Mémoire verticale non bornée}{Croissance non bornée de la profondeur de mémoire sous les retours de phase}% ADP102-0054
  \HMTTitleMapEntry{Réalisations géométriques du TRIT : régimes elliptique, parabolique et hyperbolique et double orientation temporelle}{TRIT comme état ternaire orienté : mémoire de retenue, ordre temporel et 3 régimes quadratiques}% ADP102-0055
  \HMTTitleMapEntry{APP, TRIT et TPK : support, orientation et transport}{Composition APP–TRIT–TPK : support arithmétique, orientation ternaire et transport avec mémoire}% ADP102-0056
  \HMTTitleMapEntry{TPK comme système de transport}{TPK comme dynamique de sélection, transport, mémoire, retour et déploiement}% ADP102-0057
  \HMTTitleMapEntry{Périodicité temporelle de (108) phases et lois de retour à (27), (54) et (108) itérations}{Périodicité temporelle de 108 phases et retours à 27, 54 et 108 itérations}% ADP102-0058
  \HMTTitleMapEntry{APP, TRIT et TPK : définitions directrices}{Composition générative APP→TRIT→TPK et construction de l'état enrichi}% ADP102-0059
  \HMTTitleMapEntry{Conséquences et publications du continu construit}{Connexion nonadique conjointe Γ₉ comme correspondance compatible du système inverse}% ADP102-0060
  \HMTTitleMapEntry{Conditionnement comme élagage des futurs}{Conditionnement de la mesure par restriction de l'arbre des prolongements futurs}% ADP102-0061
  \HMTTitleMapEntry{Holonomie nonadique \texorpdfstring{\(\Gamma_9\)}{Γ₉} du Topological Phase Kernel : fibre \texorpdfstring{\(\mathbb F_3^6\)}{F₃⁶}, prolongement coinductif et retour de phase avec mémoire d'état}{Holonomie nonadique Γ₉ du Topological Phase Kernel : transport de l'état enrichi, prolongement coinductif et retour 9→1 avec avancée de mémoire}% ADP102-0062
  \HMTTitleMapEntry{Connexion nonadique conjointe, compression et mémoire terminale}{Holonomie conjointe Γ₉ : compression observable T, composante orthogonale η et conservation terminale de mémoire}% ADP102-0063
  \HMTTitleMapEntry{Relations locales et holonomie}{Relations locales de transition et construction de l'holonomie nonadique Γ₉}% ADP102-0064
  \HMTTitleMapEntry{Obstructions de l'holonomie \(\Gamma_9\) : retour \(9\to1\), indices de composition et conservation de la phase, de la retenue, de la frontière, de la route et de la mémoire}{Obstructions de l'holonomie Γ₉ : retour 9→1, indices de composition et conservation de la phase, de la retenue, du bord, de la voie et de la mémoire}% ADP102-0065
  \HMTTitleMapEntry{Décomposition de la microfibre de clôture de \texorpdfstring{\(\pi\)}{π} en \texorpdfstring{\(5\)}{5} régions orientées : multiplicité \texorpdfstring{\(432+4\cdot144=1008\)}{432+4·144=1008}, axe \texorpdfstring{\(\varphi\)}{φ} et involution spéculaire \texorpdfstring{\(\pi/e\)}{π/e}}{Classification orientée de la microfibre de clôture de π : 5 régions, multiplicité 432+4·144=1008, axe φ et involution π/e}% ADP102-0066
  \HMTTitleMapEntry{Génération coinductive de \texorpdfstring{\(\pi\)}{π}, \texorpdfstring{\(e\)}{e} et \texorpdfstring{\(\varphi\)}{φ} par raffinement nonadique compatible}{Génération coinductive de π, e et φ par raffinement nonadique compatible}% ADP102-0067
  \HMTTitleMapEntry{Relèvement visible et chaîne de blocs}{Prolongement coinductif de l'émission exponentielle au moyen de blocs compatibles}% ADP102-0068
  \HMTTitleMapEntry{Signature décimale et carré positionnel}{Construction de la signature décimale au moyen du carré positionnel ternaire–décimal}% ADP102-0069
  \HMTTitleMapEntry{Retour visible et rupture de mémoire}{Retour de phase avec actualisation non triviale de la mémoire de l'état}% ADP102-0070
  \HMTTitleMapEntry{Certificat adversarial de la génération de \(e\) : unicité orbitale, récurrence factorielle, encadrements rationnels et prolongation cylindrique}{Unicité orbitale, récurrence factorielle et convergence cylindrique de la génération de e}% ADP102-0071
  \HMTTitleMapEntry{Relèvement par blocs et 1ʳᵉ frontière coinductive}{Prolongement par blocs et construction du 1er bord coinductif}% ADP102-0072
  \HMTTitleMapEntry{Du calendrier tritique au multigraphe d’incidence}{Construction du multigraphe d'incidence à partir du calendrier tritique}% ADP102-0073
  \HMTTitleMapEntry{Caractérisation, encadrements et prolongement}{Caractérisation de l'auto-échelle au moyen d'encadrements rationnels et du prolongement coinductif}% ADP102-0075
  \HMTTitleMapEntry{Sélection diagonale minimale}{Sélection diagonale minimale de cylindres compatibles et convergence à précision arbitraire}% ADP102-0076
  \HMTTitleMapEntry{Transfert de Markov compatible avec toute la limite projective}{Dynamique de Markov comme quotient modal de mémoire d'ordre 1 de l'holonomie Γ₉, compatible avec la limite inverse}% ADP102-0077
  \HMTTitleMapEntry{2e système de coordonnées : différences de pas \(3\) et \(4\)}{Système différentiel de coordonnées induit par les opérateurs de pas 3 et 4}% ADP102-0078
  \HMTTitleMapEntry{Obstructions de la prolongation affine de \(\alpha\) : télescopage, inversion de \(K\), frontière \(662/663\) et isolement des données cibles}{Obstructions du prolongement affine de α : télescopie, inversion de K, bord 662/663 et isolement des données cibles}% ADP102-0079
  \HMTTitleMapEntry{Construction de \texorpdfstring{\(\sqrt{-1}\)}{√(-1)} au moyen de \texorpdfstring{\(J^2=-I\)}{J²=-I} et coordination des géométries elliptique, parabolique et hyperbolique}{Construction de √(−1) et famille exponentielle tritique d'Euler : opérateurs J_τ et coordination des régimes elliptique, parabolique et hyperbolique}% ADP102-0080
  \HMTTitleMapEntry{Le noyau de transduction de Hensel--Witt}{Noyau de transduction Hensel–Witt : domaine, opération et préservation généalogique}% ADP102-0081
  \HMTTitleMapEntry{Noyau de transduction de Hensel–Witt : microrégions de \(\pi\), holonomie \(\Gamma_9\), registre dodécaphasique et transport hexade–octade}{Noyau de transduction Hensel–Witt : microrégions de π, holonomie Γ₉, registre dodécaphasique et transport hexade–octade}% ADP102-0082
  \HMTTitleMapEntry{Du code ternaire au système de Witt}{Transduction du code ternaire vers le système de Witt}% ADP102-0083
  \HMTTitleMapEntry{Hexade, octade et corridor \texorpdfstring{\(90/120\)}{90/120}}{Inclusion hexade–octade et transport d'incidence par les canaux 90/120}% ADP102-0084
  \HMTTitleMapEntry{Extracteur d'échelle triadique}{Opérateur triadique d'extraction d'échelle et son image rationnelle}% ADP102-0085
  \HMTTitleMapEntry{Défaut des puissances de nombres premiers et opérateur nombre}{Représentation du défaut des puissances premières au moyen de l'opérateur nombre}% ADP102-0086
  \HMTTitleMapEntry{Ablations du sélecteur et de la continuation}{Conditions d'unicité du sélecteur régional et du prolongement analytique}% ADP102-0087
  \HMTTitleMapEntry{Chaîne de génération \(\mathcal F_\pi\to b_\pi\to\rho_\pi\to D_A\to C_*^{\rm HMT}\to\Gamma_{6\to8}\) : sélection régionale, prolongation déterminantale et coordonnée orientée}{Chaîne de génération \(\mathcal F_\pi\to b_\pi\to\rho_\pi\to D_A\to C_*^{\mathrm{HMT}}\to\Gamma_{6\to8}\) : sélection régionale, prolongation déterminantale et coordonnée orientée}% ADP102-0088
  \HMTTitleMapEntry{Décomposition d'incidence hexade--octade dans les secteurs \texorpdfstring{\(q_{90}\)}{q₉₀} et \texorpdfstring{\(q_{120}\)}{q₁₂₀} : opérateur d'entrelacement \texorpdfstring{\(6\to8\)}{6→8} et structure précurseure du spectre d'aire}{Canaux d'incidence 90/120 du Topological Phase Kernel : demi-couronnes tritiques, factorisation des feuillets et réalisation hexade–octade du morphisme 6→8}% ADP102-0089
  \HMTTitleMapEntry{Géométrie dodécaphasique de mélange : antécédent angulaire CKM--CP et lecteur précurseur de l'électron}{Opérateur de mélange dodécaphasique : phase CKM–CP et construction du lecteur prédimensionnel de l'électron}% ADP102-0090
  \HMTTitleMapEntry{Morphismes archimédien et exceptionnel à partir d'un domaine enrichi commun}{Principe d'ambivalence : involution des lecteurs surfacique et volumétrique sur l'état enrichi}% ADP102-0091
  \HMTTitleMapEntry{L'opération locale sur les 729 mots}{Domaine enrichi commun et codomaines non équivalents des 2 projections}% ADP102-0092
  \HMTTitleMapEntry{Naturalité et compatibilité à toute profondeur}{Séparation conjointe des histoires par les projections archimédienne et exceptionnelle}% ADP102-0093
  \HMTTitleMapEntry{Couplage postgénératif avec la clôture opératorielle d'Euler}{Invariant d'incidence commun entre l'état enrichi et le système de Witt}% ADP102-0094
  \HMTTitleMapEntry{Globalisation naturelle des projections archimédienne et exceptionnelle sur un système projectif commun}{Décomposition de la microfibre de π en 5 régions orientées : 4 composantes périphériques et 1 axe fixe}% ADP102-0095
  \HMTTitleMapEntry{Systèmes compatibles à profondeur finie}{Prolongement en réseau de l'incidence et construction d'écrans dimensionnels}% ADP102-0096
  \HMTTitleMapEntry{Jauge de la publication exceptionnelle}{Factorisation du lecteur d'incidence à travers l'état enrichi}% ADP102-0097
  \HMTTitleMapEntry{Naturalité conjointe des publications de l'état généalogique}{Double publication de l'état enrichi : évaluation archimédienne et conservation de l'incidence}% ADP102-0098
  \HMTTitleMapEntry{Séparation conjointe de la valeur et de l'incidence exceptionnelle}{Unicité du calibre initial de la projection exceptionnelle}% ADP102-0099
  \HMTTitleMapEntry{Antécédent non commutatif commun de la double projection}{Algèbre non commutative commune et factorisation des 2 projections}% ADP102-0100
  \HMTTitleMapEntry{De l'espace des histoires à la droite réelle}{Quotient de l'espace des histoires compatibles et réalisation de la droite réelle}% ADP102-0101
  \HMTTitleMapEntry{Moyenne archimédienne et double projection du même état}{Compatibilité de l'évaluation archimédienne moyenne avec la double projection de l'état enrichi}% ADP102-0102
  \HMTTitleMapEntry{Calendrier unitaire, mémoire modulaire et mode stable}{Mode stable du calendrier unitaire sous actualisation de mémoire modulaire}% ADP102-0103
  \HMTTitleMapEntry{Publications analytiques d'histoires HMT}{Réalisations analytiques d'histoires HMT au moyen des cartes archimédienne et bidirectionnelle}% ADP102-0104
  \HMTTitleMapEntry{Conséquence pour une action positive}{Positivité de l'action induite par l'opérateur d'entrelacement dodécaphasique–Witt}% ADP102-0105
  \HMTTitleMapEntry{Du saut spectral \texorpdfstring{\(4\)}{4} au réseau de Leech : déploiement APP \texorpdfstring{\(4\to2\to6\to54\)}{4→2→6→54} et corridor Paley--Witt--Mathieu--Golay--Niemeier}{Du saut spectral 4 au réseau de Leech : déploiement APP 4→2→6→54 et corridor Paley–Witt–Mathieu–Golay–Niemeier}% ADP102-0106
  \HMTTitleMapEntry{Nature de la structure obtenue}{Classification algébrique de la structure ternaire autoduale obtenue}% ADP102-0107
  \HMTTitleMapEntry{Loi d'incidence \texorpdfstring{\(4+1\)}{4+1} sur une tétrade}{Décomposition 4+1 de l'incidence d'une tétrade en composantes périphériques et axiale}% ADP102-0108
  \HMTTitleMapEntry{Module et forme discriminants}{Décomposition de l'incidence d'une tétrade en 4 composantes périphériques et 1 composante axiale}% ADP102-0109
  \HMTTitleMapEntry{\(12\) facteurs et application de recollement}{Correspondance entre strates d'incidence et orbite libre du triplet positif}% ADP102-0110
  \HMTTitleMapEntry{Caractère et noyau de congruence}{Recollement discriminant de 12 facteurs A₂ et construction de l'extension entière}% ADP102-0111
  \HMTTitleMapEntry{Courbures mixtes d’APP}{Système ordonné d'incidences et orientation de la chambre radiale}% ADP102-0112
  \HMTTitleMapEntry{Représentation de Weyl--Heisenberg et phase centrale}{Construction du caractère radial et détermination de son noyau de congruence}% ADP102-0113
  \HMTTitleMapEntry{Statut exact de la hiérarchie \(4\to2\to6\to54\) et de l'identification réticulaire.}{Statut exact de la hiérarchie 4→2→6→54 et de l'identification de réseau.}% ADP102-0114
  \HMTTitleMapEntry{Trace graduée du générateur ternaire}{Trace graduée de l'action du générateur ternaire sur \(V^{\natural}\)}% ADP102-0115
  \HMTTitleMapEntry{Mot de support total et isométrie sans points fixes}{Isométrie sans points fixes induite par le mot ternaire de support total}% ADP102-0116
  \HMTTitleMapEntry{Construction de l'espace de Hilbert à partir de la connexion nonadique et de ses opérateurs de phase}{Réalisation projective de Bloch à partir de la structure complexe interne}% ADP102-0117
  \HMTTitleMapEntry{La branche marquée change de type}{Immersion par coins d'une phase marquée : idéal compact \(\mathcal K(\mathcal H_\infty)\) et facteur \(\mathcal B(\mathcal H_\infty)\) de type \(\mathrm I_\infty\)}% ADP102-0118
  \HMTTitleMapEntry{Reste visible et quotient de frontière}{Extension opératorielle d'APP : modes visibles et modes de mémoire}% ADP102-0119
  \HMTTitleMapEntry{Involution de dualité T}{Complétion 1000=729+243+27+1 et stratification de l'unité}% ADP102-0120
  \HMTTitleMapEntry{Complétion stratifiée de l'unité}{Décomposition 243=1+22+90+120+10 et typage de ses composantes}% ADP102-0121
  \HMTTitleMapEntry{Le bloc \texorpdfstring{\(243\)}{243} comme frontière et boule ternaire parfaite}{Dérivation des canaux de constantes à partir de la décomposition du bloc 243}% ADP102-0122
  \HMTTitleMapEntry{Décomposition interne du bloc \texorpdfstring{\(243\)}{243}}{Construction de la dimension 11 au moyen de lecteurs résiduels compatibles}% ADP102-0123
  \HMTTitleMapEntry{Dimension \texorpdfstring{\(11\)}{11} et lectures compatibles}{Clôture du registre généalogique sous l'involution de dualité T}% ADP102-0124
  \HMTTitleMapEntry{Identités vérifiables de la construction}{Identités exactes d'involution, de modularité et de conservation résiduelle}% ADP102-0125
  \HMTTitleMapEntry{Niveau, échange linéaire de masses et modularité : \(3\) contrôles de non-fusion}{Séparation par niveau, échange linéaire des masses et modularité des 3 réalisations}% ADP102-0126
  \HMTTitleMapEntry{\(2\) descendants du code ternaire étendu}{Réduction dimensionnelle 26→24→11→10 et coordination avec les écrans de théorie M}% ADP102-0127
  \HMTTitleMapEntry{La représentation icosaédrique de dimension \(12\)}{Coordination des écrans exacts avec la réalisation physique de théorie M}% ADP102-0128
  \HMTTitleMapEntry{Réduction isotrope de rang \(26\to24\)}{Représentation opératorielle d'APP et triplet spectral supersymétrique en dimension 11}% ADP102-0129
  \HMTTitleMapEntry{Coordination du corridor Moonshine avec les écrans de théorie M : graduation, réseau lorentzien et double réalisation de l'état HMT}{Théorème de compatibilité entre le corridor Moonshine et les écrans de théorie M par graduation et transport de réseau}% ADP102-0130
  \HMTTitleMapEntry{Antécédent commun des lecteurs : raffinement spectral et limite de dilatation conjointe}{Raffinement spectral commun et limite de la dilatation conjointe des lecteurs matriciels}% ADP102-0131
  \HMTTitleMapEntry{Sous-structures, objet global et projections}{Construction de 4 carrés latins quantiques cycliques dans une classe opératorielle élémentaire}% ADP102-0132
  \HMTTitleMapEntry{Correspondances directes et inverses entre les formalismes}{Réalisation de 3 contextes opératoriels dans un carré quantique (9,3)}% ADP102-0133
  \HMTTitleMapEntry{Extension de la bigraduation à 3 indices}{Représentation de l'atlas généalogique complet dans un carré quantique (12,3)}% ADP102-0134
  \HMTTitleMapEntry{4 carrés latins quantiques cycliques de classe opératorielle facile}{Classification cyclique de 4 carrés latins quantiques dans une classe opératorielle élémentaire}% ADP102-0135
  \HMTTitleMapEntry{Lemme cyclique des carrés magiques quantiques associés à une mesure opératorielle positive}{Bigraduation de la famille généalogico–matricielle}% ADP102-0136
  \HMTTitleMapEntry{Cônes matriciels, images et extensions complètement positives}{Représentation matricielle de la composition généalogique APP–TRIT–TPK}% ADP102-0137
  \HMTTitleMapEntry{Construction explicite de cubes magiques quantiques et compatibilité de leurs marges}{TPK comme catégorie de voies avec sélection, transport, mémoire et retour}% ADP102-0138
  \HMTTitleMapEntry{Combinaisons magiques d'applications CP}{Compositions complètement positives compatibles avec la bigraduation généalogico–matricielle}% ADP102-0139
  \HMTTitleMapEntry{Hiérarchie spectrale \(4\to2\to6\to54\) du bloc central d'APP}{Réalisation métrologique ultérieure de la double projection de l'état enrichi}% ADP102-0140
  \HMTTitleMapEntry{Raffinement nonadique par espérance conditionnelle}{Décomposition symplectique 12→4×3 des coordonnées de phase}% ADP102-0141
  \HMTTitleMapEntry{Taxonomie fonctionnelle de \(114\) expressions numériques et de \(106\) lectures opératorielles}{Classification opératorielle de 114 expressions numériques par domaine, générateur, invariant et codomaine}% ADP102-0142
  \HMTTitleMapEntry{Transfert de phase sur le multigraphe des blocs complets}{Microfibre de clôture de π : 5 régions orientées, multiplicités \(432+4\cdot144=1008\) et projection commune \(w_\pi\)}% ADP102-0143
  \HMTTitleMapEntry{\(5\) cycles orientés de clôture de la microfibre}{Publication de mémoire réduite d'un bloc nonadique complet au moyen du multigraphe de transfert de phase}% ADP102-0144
  \HMTTitleMapEntry{Opérateur d’incidence et demi-tour orienté}{Holonomies orientées de clôture des 5 régions de π}% ADP102-0145
  \HMTTitleMapEntry{Composition rationnelle et compensateur \texorpdfstring{\(239\)}{239}}{Détermination de la pente 1/5 par la multiplicité de la microfibre régionale de π}% ADP102-0146
  \HMTTitleMapEntry{Identité gaussienne et \(1^{\mathrm a}\) demi-tour orienté}{Composition rationnelle du demi-tour et détermination du terme compensateur 239}% ADP102-0147
  \HMTTitleMapEntry{Prolongement coinductif de la section de \(\pi\) par le système inverse des histoires survivantes}{Filtration de F₃⁶ par admissibilité régionale et sélection de sous-cylindres compatibles}% ADP102-0148
  \HMTTitleMapEntry{Système inverse des histoires compatibles}{Unicité de la section diagonale minimale dans le système d'encadrements rationnels}% ADP102-0149
  \HMTTitleMapEntry{Signature positionnelle décimale et carré de compatibilité}{Prolongement de la section exponentielle au moyen de la chaîne de blocs visibles}% ADP102-0150
  \HMTTitleMapEntry{Retour de phase avec changement d’état et conservation de la mémoire parcourue}{Construction de la signature décimale et du carré positionnel de la section exponentielle}% ADP102-0151
  \HMTTitleMapEntry{Semi-groupe factoriel du caractère exponentiel : récurrence \(a_n=1/n!\)}{Retour de phase avec inégalité des états enrichis : \(g(K+9)=g(K)\) et \(B_{K+9}\ne B_K\)}% ADP102-0152
  \HMTTitleMapEntry{Certificat interne d’unicité orbitale, de récurrence factorielle, de convergence et de prolongement cylindrique}{Prolongement coinductif de e au moyen de l'holonomie nonadique Γ₉ du TPK}% ADP102-0153
  \HMTTitleMapEntry{Demi-droite positive invariante et unicité de l’auto-échelle}{Construction du multigraphe d'incidence modale à partir du calendrier tritique}% ADP102-0154
  \HMTTitleMapEntry{Certificat interne d’unicité, de compatibilité et de convergence de l’auto-échelle de Perron}{Caractérisation de φ par le rayon de Perron, l'emboîtement des encadrements et le prolongement coinductif}% ADP102-0156
  \HMTTitleMapEntry{Dérivations internes de la constante de structure fine \(\alpha\) : clôture entière dodécaphasique réversible et racine analytique unique du noyau \(E_{21/22}\)}{Clôture dodécaphasique de \(\alpha_{\mathrm{HMT}}\) et caractérisation analytique au moyen de \(E_{21/22}\) : génération réversible, racine simple et prolongement nonadique d'ordre 9}% ADP102-0157
  \HMTTitleMapEntry{État terminal commun et 2 lectures internes coordonnées}{Évaluations terminale et dodécaphasique de l'état enrichi commun : blocs \(B_{\mathrm{term}}\) et carte signée \(U_{12}^{\mathrm{sgn}}\)}% ADP102-0158
  \HMTTitleMapEntry{Voie A : clôture entière dodécaphasique}{Génération réversible de \(\alpha_{\mathrm{HMT}}\) au moyen de la clôture dodécaphasique entière}% ADP102-0159
  \HMTTitleMapEntry{Domaine Hensel--Witt, applications de Teichmüller et cible rationnelle \(-1/12\)}{Opérateur Hensel–Witt de transduction entre microrégions de π, holonomie Γ₉, registre dodécaphasique et incidence hexade–octade}% ADP102-0160
  \HMTTitleMapEntry{Application de Teichmüller du code ternaire à l’anneau de Witt}{Construction du système de Witt à partir du codage ternaire compatible}% ADP102-0161
  \HMTTitleMapEntry{Incidence hexade--octade et canaux modulaires \(90/120\)}{Réalisation hexade–octade des canaux 90/120 précédemment dérivés de l'incidence tritique du TPK}% ADP102-0162
  \HMTTitleMapEntry{4 routes compatibles vers \texorpdfstring{\(-1/12\)}{-1/12}}{4 réalisations opératorielles de −1/12 et séparation de leurs domaines}% ADP102-0163
  \HMTTitleMapEntry{Dérivation déterminantale de l’échelle d’action : forme normale de Smith, conoyau d’ordre \(16\) et valuation \(10^{-34}\)}{Dérivation déterminantale de l'échelle absolue d'action : forme normale de Smith, conoyau d'ordre 16 et valuation 10⁻³⁴}% ADP102-0164
  \HMTTitleMapEntry{Mesure de quotient et caractère positif de l’action : hilbertisation canonique, défaut régional--dodécaphasique et loi fonctorielle sur les chemins}{Construction de la mesure positive d'action sur le quotient des émissions : hilbertisation, défaut régional–dodécaphasique et fonctorialité des chemins}% ADP102-0165
  \HMTTitleMapEntry{Dérivation autoduale de la constante de Planck : périodicité temporelle de \(108\) phases, action et longueur de Compton}{Dérivation autoduale de la constante de Planck et de sa droite d'action : périodicité temporelle de 108 phases, fermeture de Compton et covariance bidirectionnelle}% ADP102-0166
  \HMTTitleMapEntry{Équivalence des normalisations sur l’identité autoduale}{Commutativité des normalisations angulaire, temporelle et gravitationnelle de l'identité autoduale}% ADP102-0167
  \HMTTitleMapEntry{Opérateur angulaire prédimensionnel : décomposition en \(27\) secteurs, rayon spectral et composition effective de la double projection}{Opérateur de phase antérieur à la transduction dimensionnelle : décomposition en 27 secteurs, rayon spectral et factorisation de la double projection}% ADP102-0168
  \HMTTitleMapEntry{Caractères quadratiques et constantes spectrales : Catalan, Apéry et tour d’Euler--Stieltjes}{Caractères quadratiques et constantes spectrales : Catalan, Apéry, Euler–Mascheroni et hiérarchie de Stieltjes}% ADP102-0169
  \HMTTitleMapEntry{Criticité unimodale de la famille dodécaphasique : profil analytique exact, approximants de Feigenbaum et réduction de l’universalité au certificat hyperbolique}{Criticité unimodale de la famille dodécaphasique : profil analytique exact, approximants de Feigenbaum et certificat hyperbolique de renormalisation}% ADP102-0170
  \HMTTitleMapEntry{Arithmétique du mode \(30\) : identités de Rogers--Ramanujan, normalisateur \(899\) et relèvement nonadique de Pell}{Arithmétique de la famille modulaire dodécaphasique de période 30 : identités de Rogers–Ramanujan, normalisateur 899 et relèvement nonadique de Pell}% ADP102-0171
  \HMTTitleMapEntry{Théorème directeur du réseau radical et de la loi puissance--lacune}{Classification du réseau radical 3→6→9 et dérivation de la loi puissance–lacune}% ADP102-0172
  \HMTTitleMapEntry{Orbite cyclique \(369\to693\to936\to369\)}{Orbite cyclique 369→693→936→369}% ADP102-0173
  \HMTTitleMapEntry{Noyau \(7227\), décomposition spectrale et lien avec la croix radicale}{Détermination du noyau arithmétique 7227 par incidence avec la croix radicale d'APP}% ADP102-0174
  \HMTTitleMapEntry{Encodage mécanique des lacunes par les puissances nonadiques}{Codage sturmien des lacunes induit par les puissances nonadiques}% ADP102-0175
  \HMTTitleMapEntry{Racine interne de \(-1\), structure exponentielle et géométries elliptique, parabolique et hyperbolique}{Famille exponentielle triadique \(J_\tau^2=-\tau I\) : réalisations elliptique, parabolique et hyperbolique de l'identité d'Euler}% ADP102-0176
  \HMTTitleMapEntry{Réalisation continue du type quadratique induit}{Extension continue de l'opérateur quadratique commun et conservation de son type spectral}% ADP102-0177
  \HMTTitleMapEntry{Support opératoriel commun de \(3\) lectures géométriques}{Conjugaison des réalisations arithmétique, angulaire et torsionnelle au moyen de l'opérateur quadratique commun}% ADP102-0178
  \HMTTitleMapEntry{Chaîne constitutive du secteur gravitationnel : échelle chronologique \(108\), réduction tensorielle, holonomie et évaluation métrologique de \(G\)}{Composition causale des invariants de phase, d'action et de longueur dans la dérivation de G}% ADP102-0179
  \HMTTitleMapEntry{Cardinalités d’incidence 90 et 120}{Dérivation tritique des cardinalités 90/120 et réalisation ultérieure dans l'incidence hexade–octade}% ADP102-0180
  \HMTTitleMapEntry{Séparation causale entre le spectre d’aire, la mémoire modulaire et la réalisation Cartan--Holst}{Distinction typée entre Barbero–aire, aire–hamiltonien modulaire et réalisation Cartan–Holst}% ADP102-0182
  \HMTTitleMapEntry{Transducteur de Hensel non filtré et décalage complet sur les fibres ternaires}{Transducteur entier de Hensel comme opérateur de décalage complet}% ADP102-0183
  \HMTTitleMapEntry{Surjectivité de la carte archimédienne non filtrée sur \(\mathbb R\)}{Couverture exacte de la carte réelle antécédente au moyen de sections compatibles}% ADP102-0184
  \HMTTitleMapEntry{2 axes de classification : valeur visible et généalogie de transport}{Classification bifactorielle du nombre par valeur archimédienne et généalogie de transport}% ADP102-0185
  \HMTTitleMapEntry{Position de \texorpdfstring{\(\pi,e,\varphi\) et \(\alpha\)}{pi, e, phi et alpha} dans la classification}{Position de π,e,φ et α dans la classification}% ADP102-0186
  \HMTTitleMapEntry{Carte de Hensel non filtrée, survie et frontière globale}{Carte antécédente, survie coinductive et construction du bord global}% ADP102-0187
  \HMTTitleMapEntry{La famille opératorielle d'Euler comme grammaire de clôture}{Factorisation de la fermeture généalogique du nombre au moyen du système opératoriel d'Euler}% ADP102-0188
  \HMTTitleMapEntry{Dépendances mathématiques.}{Dépendance causale entre section inverse, évaluation archimédienne et conservation exceptionnelle de l'incidence}% ADP102-0189
  \HMTTitleMapEntry{\(4\) exigences de choix dans le système projectif}{Conditions de choix nécessaires au prolongement global des sections compatibles}% ADP102-0190
  \HMTTitleMapEntry{Décidabilité effective de chaque fenêtre finie de multiplicités projectives}{Décidabilité algorithmique de chaque niveau fini du système de préfixes}% ADP102-0191
  \HMTTitleMapEntry{Relation exacte avec les \(9\) phases du Topological Phase Kernel}{Restriction de l'obstruction uniforme au système temporel nonadique à 9 phases}% ADP102-0192
  \HMTTitleMapEntry{De la cohérence constructible à l'indépendance par forcing}{Cohérence relative dans L et indépendance par forcing}% ADP102-0193
  \HMTTitleMapEntry{Optimisation minimax sur des fenêtres finies de préfixes compatibles}{Théorème minimax à chaque niveau fini de l'ensemble partiellement ordonné des préfixes compatibles}% ADP102-0194
  \HMTTitleMapEntry{La chaîne que le pont reçoit démontrée}{Antécédents théorématiques du foncteur de transduction généalogico–métrologique}% ADP102-0195
  \HMTTitleMapEntry{Dérivation bidirectionnelle de l'exposant électronique à partir du parcours enrichi}{Bande de réalisation de l'électron et séparation entre masse basale et clôture bêta}% ADP102-0196
  \HMTTitleMapEntry{\(2\) branches causales et leur confluence}{Confluence des branches déterminantale et dodécaphasique de l'échelle d'action}% ADP102-0197
  \HMTTitleMapEntry{Carte métrologique ultérieure et évaluation dimensionnelle des invariants engendrés}{Systèmes de transduction entre nombre généalogique et grandeur physique}% ADP102-0198
  \HMTTitleMapEntry{Principe d'ambivalence : involution des lecteurs surfacique et volumétrique, équivalence locale et orientation temporelle conjuguée}{Principe d'ambivalence en mécanique dimensionnelle : involution surfacique–volumétrique, systèmes auto-inertiels et orientation temporelle conjuguée}% ADP102-0199
  \HMTTitleMapEntry{Conditions d'une évaluation prospective}{Conditions d'évaluation prospective du principe d'ambivalence sur l'état enrichi}% ADP102-0200
  \HMTTitleMapEntry{Cristal temporel fini HMT : forçage, sous-harmonique et stabilité spectrale}{Cristal temporel fini HMT : hamiltonien périodiquement forcé, réponse sous-harmonique et stabilité spectrale}% ADP102-0201
  \HMTTitleMapEntry{Forçage construit à partir de l’horloge à 108 états}{Construction du forçage périodique à partir du système temporel à 108 états}% ADP102-0202
  \HMTTitleMapEntry{Stabilité spectrale finie}{Jet nilpotent de l'évolution triadique et prolongement compatible}% ADP102-0203
  \HMTTitleMapEntry{Théorème de confluence HMT–MD : forme électrogravitationnelle, système d’Einstein–Schrödinger et coordination de l’action, de la matière, des champs, de la géométrie et de l’information}{Théorème de confluence HMT–MD : système Einstein–Schrödinger et coordination électrogravitationnelle de l'action, de la matière, des champs, de la géométrie et de l'information}% ADP102-0204
  \HMTTitleMapEntry{Indépendance horizontale et compatibilité verticale}{Naturalité des réalisations et conservation verticale des invariants de l'état enrichi}% ADP102-0205
  \HMTTitleMapEntry{1re confluence : couplage, vide et secteur électrofaible}{Couplage constitutif entre le vide et le secteur électrofaible}% ADP102-0206
  \HMTTitleMapEntry{2e confluence : action, gravité et forme électrogravitationnelle}{Réciprocité entre action et gravité dans la forme électrogravitationnelle}% ADP102-0207
  \HMTTitleMapEntry{3e confluence : géométrie, information et mélange}{Confluence entre géométrie, information et mélange}% ADP102-0208
  \HMTTitleMapEntry{4e confluence : écran, torsion et mémoire hélicoïdale}{Transduction de la mémoire hélicoïdale dans la réalisation torsionnelle}% ADP102-0209
  \HMTTitleMapEntry{Du déterminant électrofaible à l’opérateur constitutif positif du vide et à sa complétion spectrale \(3\to4\)}{Du déterminant électrofaible à l'opérateur constitutif positif du vide et à sa complétion spectrale 3→4}% ADP102-0210
  \HMTTitleMapEntry{Valeur publiée et mémoire de la fibre}{Désintégration spectrale de la mesure des trajectoires et conditionnement des futurs}% ADP102-0211
  \HMTTitleMapEntry{Fondements algébriques de la couture de phase et de l'automesure}{État conjoint observateur–observable, instrument réversible et automesure}% ADP102-0212
  \HMTTitleMapEntry{Information paire et transport géométrique impair dans la bicouche}{Décomposition paire–impaire de l'état bicouche : information interfeuillet et transport géométrique orienté}% ADP102-0214
  \HMTTitleMapEntry{\texorpdfstring{\(3\)}{3} réalisations de la mémoire généalogique : transport dodécaphasique, fonction de Green non locale avec générateur local et expansion cosmologique}{Réalisations dynamique, propagative et cosmologique de la mémoire généalogique : transport dodécaphasique, fonction de Green et expansion spatiale}% ADP102-0215
  \HMTTitleMapEntry{Chaîne constitutive du secteur gravitationnel : échelle chronologique \texorpdfstring{\(108\)}{108}, réduction tensorielle, holonomie et évaluation métrologique de \texorpdfstring{\(G\)}{G}}{Descente icosaédrique A₅/C₅ et sélection du porteur gravitationnel}% ADP102-0216
  \HMTTitleMapEntry{Entrelaceur avec les tenseurs symétriques sans trace et réduction transverse \texorpdfstring{\(12\to5\to5\to2\)}{12→5→5→2}}{Opérateur d'entrelacement avec les tenseurs symétriques sans trace et réduction transverse 12→5→5→2}% ADP102-0217
  \HMTTitleMapEntry{Critères structurels du sélecteur gravitationnel : transducteur dimensionnel, holonomie enrichie et séparation entre sélecteurs torsionnel et circulaire}{Sélecteur holonomique gravitationnel : transduction dimensionnelle, naturalité, raffinement et séparation des régimes torsionnel et circulaire}% ADP102-0218
  \HMTTitleMapEntry{Carte décimale déclarée et comparaison métrologique}{Reconnaissance métrologique ultérieure de G et indépendance causale de la dérivation}% ADP102-0219
  \HMTTitleMapEntry{Tétrade, connexion et \(2\) défauts}{Tétrade, connexion de spin et torsion de la réalisation Cartan–Holst}% ADP102-0220
  \HMTTitleMapEntry{Action du 1er ordre}{Action de Palatini–Holst et équations variationnelles}% ADP102-0221
  \HMTTitleMapEntry{Confluence horizon–aire–information–entropie–action}{Théorème de confluence entre énergie d'horizon, aire, information, entropie et action}% ADP102-0222
  \HMTTitleMapEntry{\(2\) orientations d’une même propagation}{Propagateurs avancé et retardé sur des voies orientées}% ADP102-0223
  \HMTTitleMapEntry{Orbite elliptique, action et \(3\) invariants J}{Invariants de l'involution temporelle et conjugaison particule–antiparticule}% ADP102-0224
  \HMTTitleMapEntry{Cosmologie de périodicité \texorpdfstring{\(108\)}{108} : inversion de feuillet, horizon, Big Bang, rebond torsionnel et secteur sombre}{Cosmologie torsionnelle de périodicité 108 : expansion, retour avec mémoire et rebond d'Einstein–Cartan}% ADP102-0225
  \HMTTitleMapEntry{Séparation typologique des \(2\) réalisations cosmologiques}{Séparation entre accélération de de Sitter et rebond torsionnel d'Einstein–Cartan}% ADP102-0226
  \HMTTitleMapEntry{\(3\) feuilletages exacts de de Sitter}{Feuilletages cosmologiques compatibles avec l'inversion de feuillet}% ADP102-0227
  \HMTTitleMapEntry{Univers à frontière}{État de bord de l'horizon cosmologique et mémoire de retour}% ADP102-0228
  \HMTTitleMapEntry{Données spectrales d'une réalisation supersymétrique : \((\mathcal H_{\rm phys},H,\mathcal Q)\), supercharge et compactification en dimension \(11\)}{Réalisation de théorie M par compactification, cordes, branes et dimension 11}% ADP102-0229
  \HMTTitleMapEntry{Loi générale des masses sur l'atlas \texorpdfstring{\(81/56/13/324\)}{81/56/13/324} : trajectoires, registre, signature, caractère, tours \texorpdfstring{\(90/120\)}{90/120} et secteur nul}{Loi générale des masses sur l'atlas 81/56/13/324 : trajectoires, registre, signature, caractère, semi-groupe harmonique ⟨90,120⟩ et secteur nul}% ADP102-0230
  \HMTTitleMapEntry{Algorithme universel de réalisation matérielle et égalisateur spectral des lecteurs \(K_D\) et \(K_\Omega\)}{Construction fonctorielle des familles de masse sur l'atlas}% ADP102-0231
  \HMTTitleMapEntry{Bifurcation nulle et positivité}{Bifurcation du secteur nul antérieure à la branche positive de masse}% ADP102-0232
  \HMTTitleMapEntry{Accès au domaine exhaustif}{Construction exhaustive de l'atlas 81/56/13/324 et conservation du secteur nul}% ADP102-0233
  \HMTTitleMapEntry{Domaine interne et problème de réalisation}{Théorème d'exhaustivité du domaine matériel et stabilité du secteur nul sous raffinement}% ADP102-0234
  \HMTTitleMapEntry{Produit fibré entre descripteurs et routes}{Produit fibré entre invariants d'espèce et voies généalogiques}% ADP102-0235
  \HMTTitleMapEntry{Lecteurs bidirectionnels \texorpdfstring{\(K_D\)}{K D} et \texorpdfstring{\(K_\Omega\)}{K Omega} : profils, coïncidences, tours \texorpdfstring{\(90/120\)}{90/120} et opérateurs par multisection}{Lecteurs bidirectionnels \(K_D\) et \(K_\Omega\) : profils, coïncidences, semi-groupe harmonique \(\langle90,120\rangle\) et opérateurs par multisection}% ADP102-0236
  \HMTTitleMapEntry{Lecteurs \(K_D\) et \(K_\Omega\) : profils dodécaphasiques, domaine des parcours et codomaine spectral}{Lecteurs \(K_D\) et \(K_\Omega\) : profils dodécaphasiques, domaine des parcours et codomaine spectral}% ADP102-0237
  \HMTTitleMapEntry{Réalisations spectrales, scalaires et nulles de la Loi générale des masses : familles matérielles, transport CKM--PMNS et secteurs photonique et gluonique}{Réalisations spectrales, scalaires et nulles de la Loi générale des masses : familles matérielles, mélange quark CKM, mélange neutrinique PMNS et secteurs nuls photonique et gluonique}% ADP102-0238
  \HMTTitleMapEntry{Couleur, saveurs de quarks et mélange CKM}{Opérateurs de masse par fibre, orbite et multisection de l'atlas}% ADP102-0239
  \HMTTitleMapEntry{États composés et loi de liaison}{Correspondance par domaine et codomaine entre l'inventaire matériel et les réalisations de l'atlas}% ADP102-0240
  \HMTTitleMapEntry{Classificateur des espèces matérielles par fibre, représentation, charge, spin, saveur et génération}{Classification des espèces matérielles par fibre, représentation, charge, spin, saveur et génération}% ADP102-0241
  \HMTTitleMapEntry{Typage probatoire de la persistance, de l'atlas et des familles.}{Conditions de persistance et correspondance entre atlas, familles et représentations}% ADP102-0242
  \HMTTitleMapEntry{Ontologie matérielle, charge, spin, statistique, couleur et saveur}{Définitions structurelles de particule, masse, charge, spin, statistique, couleur, saveur et génération}% ADP102-0243
  \HMTTitleMapEntry{Théorème de clôture typée de la réalisation massique}{Théorème de clôture de la réalisation matérielle au moyen de sections persistantes et de représentations finies}% ADP102-0244
  \HMTTitleMapEntry{Confrontation métrologique et archive exhaustive des réalisations}{Reconnaissance métrologique ultérieure et inventaire exhaustif des réalisations matérielles}% ADP102-0245
  \HMTTitleMapEntry{Architecture naturelle commune des réalisations terminales du continu généalogique}{Théorème de factorisation conjointe des \(5\) réalisations terminales à travers la structure discrète du continu}% ADP102-0246
  \HMTTitleMapEntry{Publication terminale du continu unique par \(5\) lecteurs non permutables}{Publication terminale au moyen de 5 lecteurs non permutables et facteur commun de l'état enrichi}% ADP102-0247
  \HMTTitleMapEntry{Fermetures dynamiques du continu : transport solénoïdal et quotient de jauge}{Dynamique et courbure du continu généalogique antérieures aux réalisations classiques}% ADP102-0248
  \HMTTitleMapEntry{Dynamique et courbure du continu généalogique}{Factorisation conjointe des réalisations terminales à partir du même état du continu}% ADP102-0249
  \HMTTitleMapEntry{Fermetures algébriques du continu et limites de la reconstruction extensionnelle}{Factorisation commune, cohérence extensionnelle et limites de la reconstruction terminale}% ADP102-0250
  \HMTTitleMapEntry{L’exemple des ordinaux de von Neumann}{Représentation ordinale de von Neumann comme modèle de cohérence extensionnelle}% ADP102-0251
  \HMTTitleMapEntry{Domaine, applications et compatibilités de l’objet terminal commun}{Vue interne \(G_T^{\mathrm{int}}\) du continu et application \(\operatorname{App}_{T,d_T}\) : conservation de la fibre, de la retenue, de la mémoire et du terminal}% ADP102-0252
  \HMTTitleMapEntry{Déploiement interne des réalisations corrélatives avant leur identification aux problèmes classiques}{Émergence corrélative des 5 actions terminales sur le même état du continu}% ADP102-0253
  \HMTTitleMapEntry{\(5\) lecteurs terminaux non permutables du continu généalogique}{Théorème de non-permutabilité des 5 lecteurs terminaux et classification de leurs codomaines}% ADP102-0254
  \HMTTitleMapEntry{Obstructions à la projection ablative du continu : perte de phase, d'orientation ou de mémoire généalogique}{Projecteurs positifs de la forme de Guinand–Weil et localisation spectrale}% ADP102-0255
  \HMTTitleMapEntry{Fixation de la droite critique par l’involution fonctionnelle \(s\mapsto1-\bar s\)}{Fixation de la droite critique par l'involution fonctionnelle \texorpdfstring{\(s\mapsto 1-\bar{s}\)}{s -> 1 - conjugué(s)}}% ADP102-0256
  \HMTTitleMapEntry{Distinction typée entre les zéros de \(\zeta\), de \(\xi\) et du déterminant spectral}{Séparation typée des zéros de ζ, ξ et du déterminant spectral}% ADP102-0257
  \HMTTitleMapEntry{Représentant positif de la classe résiduelle nulle : \(n\equiv0\pmod 9\) et \(\rho_9(n)=9\)}{Représentant positif de la classe résiduelle nulle : \(n\equiv0\pmod 9\) et \(\rho_9(n)=9\)}% ADP102-0258
  \HMTTitleMapEntry{Isométrie de transport tordu \(\widetilde U_k\) sur les états cylindriques, les poids projectifs et la mémoire de Weil}{Isométrie de transport tordu \(\widetilde U_k\) sur les états cylindriques, les poids projectifs et la mémoire de Weil}% ADP102-0259
  \HMTTitleMapEntry{Invariant gradué du corridor spectral \(4\to2\to6\to54\)}{Invariant gradué du corridor spectral 4→2→6→54}% ADP102-0260
  \HMTTitleMapEntry{Trace couplée du spectre de Weil et de la graduation Moonshine}{Construction indépendante du champ spectral sur le double cercle}% ADP102-0261
  \HMTTitleMapEntry{Domaine hydrodynamique et publication de Galerkin}{Domaine solénoïdal et conservation de la divergence nulle}% ADP102-0262
  \HMTTitleMapEntry{Donnée unique de départ et noyaux denses}{Données initiales, compatibilité de bord et propagation régulière}% ADP102-0263
  \HMTTitleMapEntry{Réalisation finie de l'inégalité de Kalman–Yakubovich–Popov et fonction auxiliaire dans la clôture énergétique}{Inégalité KYP et contrôle coercif de l'évolution solénoïdale}% ADP102-0264
  \HMTTitleMapEntry{Extraction matérielle de la borne de Sobolev}{Estimation de Sobolev et exclusion de la concentration singulière}% ADP102-0265
  \HMTTitleMapEntry{La fibre d'incidence est \(1\) coordonnée physique de l'état complet}{Courbure de fibre et coercivité de l'action de jauge}% ADP102-0266
  \HMTTitleMapEntry{Vide simple et témoin qui atteint l'écart spectral}{Unicité du vide invariant de jauge et saturation de l'écart \(3\kappa_{\mathrm{av}}\) par le vecteur propre \(W_c\)}% ADP102-0267
  \HMTTitleMapEntry{\(4\) reconstructions d'Osterwalder–Schrader d'une dynamique commune}{Positivité d'Osterwalder–Schrader et reconstruction quantique}% ADP102-0268
  \HMTTitleMapEntry{Théorie de jauge euclidienne et gap de Yang--Mills}{Écart spectral du secteur de jauge et limite projective}% ADP102-0269
  \HMTTitleMapEntry{Présentation générateur par générateur de la nouvelle carte}{Générateurs cohomologiques et projecteurs compatibles}% ADP102-0270
  \HMTTitleMapEntry{Fibre initiale, Mittag--Leffler et publication affirmative}{Lefschetz–Murre et descente algébrique des classes de Hodge}% ADP102-0271
  \HMTTitleMapEntry{Bocksteins de toutes les pages}{Homomorphismes de Bockstein de la suite spectrale : compatibilité entre pages et détermination de l'ordre en T du complexe universel}% ADP102-0272
  \HMTTitleMapEntry{Table complète Kato--FERS avant le déterminant}{Déterminant de Kato–FERS, rang et terme principal de BSD}% ADP102-0273
  \HMTTitleMapEntry{Distinction entre structure et application}{Vue interne \(G_T\) et lecteur terminal \(P_T\) : indépendance logique entre existence structurelle et réalisation classique}% ADP102-0274
  \HMTTitleMapEntry{Obstructions communes aux \(5\) expressions terminales : perte de domaine, de naturalité, de certificat ou de limite compatible}{Obstructions communes des lecteurs terminaux : perte de domaine, de naturalité ou de limite compatible}% ADP102-0275
  \HMTTitleMapEntry{Formalisation assistée du noyau algébrique}{Formalisation vérifiable du noyau algébrique et certification d'identités finies}% ADP102-0276
  \HMTTitleMapEntry{Correspondance entre énoncés, preuves, certificats et propriétaires scientifiques}{Correspondance entre énoncés, démonstrations, certificats et sources scientifiques primaires}% ADP102-0277
  \HMTTitleMapEntry{Conventions de typage et de représentation des données}{Conventions mathématiques et représentation des données}% ADP102-0278
  \HMTTitleMapEntry{Hétérogénéité typée et couverture du catalogue}{Hétérogénéité des domaines physiques et exhaustivité du catalogue}% ADP102-0279
  \HMTTitleMapEntry{Fiches individuelles des chemins orientés}{Déploiement intégral des 324 voies orientées : correspondance entre 53 champs canoniques, 57 champs enrichis et 49 invariants publiés par section}% ADP102-0280
  \HMTTitleMapEntry{Inventaire universel typé de 471 masses et 384 largeurs}{Inventaire exhaustif de 471 masses et 384 largeurs : identification, provenance et correspondance avec l'atlas}% ADP102-0281
  \HMTTitleMapEntry{Source, déduplication et limites de type}{Provenance des données, unicité des enregistrements et domaines de validité}% ADP102-0282
  \HMTTitleMapEntry{Code d’objet et lecteur HMT--MD}{Identification des objets et application du lecteur HMT--MD}% ADP102-0283
  % Correspondencias espirales: sustituciones científicas unívocas aprobadas
  % por cotejo focal. Se modifican sólo los rótulos; los cuerpos permanecen
  % íntegros en sus residencias distribuidas.
  \HMTTitleMapEntry{TPK : sélection, transport, actualisation et retour}{Composition dynamique du Topological Phase Kernel sur l'état enrichi : sélection tritique, transport cardinal, actualisation de mémoire et retour typé}% ADP102-0284
  \HMTTitleMapEntry{La transition typée}{Définition et factorisation de l'opérateur de transition enrichie \(U_t\)}% ADP102-0285
  \HMTTitleMapEntry{Composition de la retenue}{Loi de cocycle bimodulaire de la retenue sous concaténation des voies}% ADP102-0286
  \HMTTitleMapEntry{Calendrier dodécaphasique et non-scindement}{Extension cyclique non scindée \(0\to C_{12}\to C_{108}\to C_9\to0\) et chronologie interne du TPK}% ADP102-0287
  \HMTTitleMapEntry{Hélice nonadique de phase et de profondeur}{Revêtement hélicoïdal de la phase nonadique et coordonnée entière de profondeur de mémoire}% ADP102-0288
  \HMTTitleMapEntry{Immersion géométrique normalisée}{Immersion \(H_9\) du quotient--résidu nonadique dans \(\mathbb R^3\) et calcul des invariants de Frenet}% ADP102-0289
  \HMTTitleMapEntry{Cercles emboîtés et cylindres}{Sections circulaires et feuilletage cylindrique du revêtement phase--mémoire}% ADP102-0290
  \HMTTitleMapEntry{Feuillets orbitaux et couture de l'espace de suspension}{Orbites du produit croisé et relation d'identification de l'espace de suspension}% ADP102-0291
  \HMTTitleMapEntry{Cas stationnaire affine}{Itération de l'endomorphisme affine stationnaire et classification de ses points fixes}% ADP102-0292
  \HMTTitleMapEntry{Inversion des voies}{Anti-involution du sous-groupoïde des voies réversibles et conjugaison de l'holonomie affine}% ADP102-0293
  \HMTTitleMapEntry{Miroir des constantes et axe fixe}{Involution \(\pi/e\) de la fibre des constantes, invariant \(u^\pi+u^e\) et axe fixe \(\varphi\)}% ADP102-0294
  \HMTTitleMapEntry{La géométrie visible comme publication de l'état enrichi}{Évaluation archimédienne de l'état enrichi et défaut de fidélité généalogique}% ADP102-0295
  \HMTTitleMapEntry{L'inversion de l'ordre explicatif}{Factorisation causale de l'histoire enrichie à l'évaluation archimédienne}% ADP102-0296
  \HMTTitleMapEntry{État, histoire et publication}{Espace des états enrichis, catégorie des histoires TPK et évaluation archimédienne}% ADP102-0297
  \HMTTitleMapEntry{Cercle intrinsèque et cercle projeté}{Distinction entre orbite radiale intrinsèque et projection circulaire d'une trajectoire hélicoïdale}% ADP102-0298
  \HMTTitleMapEntry{La cellule nonadique et la naissance bidimensionnelle}{Construction du support nonadique bidimensionnel \(X_9=(\mathbb Z/9\mathbb Z)^2\)}% ADP102-0299
  \HMTTitleMapEntry{Antécédents APP de la classification radiale}{Invariants APP qui induisent le caractère radial accumulé}% ADP102-0300
  \HMTTitleMapEntry{État ternaire relevé}{État TRIT enrichi par résidu et retenue : multiplicité des représentants du seuil}% ADP102-0301
  \HMTTitleMapEntry{Algèbres quadratiques}{Famille \(\mathcal A_\tau=\mathbb R[J]/(J^2+\tau I)\) et classification elliptique, parabolique et hyperbolique}% ADP102-0302
  \HMTTitleMapEntry{Double projection}{Factorisation conjointe des publications archimédienne et exceptionnelle de l'état enrichi}% ADP102-0303
  \HMTTitleMapEntry{Représentation locale d'une transition}{Représentation affine \(e\mapsto h_e\) des générateurs de la catégorie des voies}% ADP102-0304
  \HMTTitleMapEntry{Holonomies préfixes}{Produits affines ordonnés \texorpdfstring{\(H_j=h_{e_j}\cdots h_{e_1}\)}{H_j = h_ej ... h_e1} et conservation des préfixes}% ADP102-0305
  \HMTTitleMapEntry{Catégorie relevée et caractères élémentaires}{Catégorie des voies enrichies par 2 orientations TRIT et définition des caractères \(p\) et \(a\)}% ADP102-0306
  \HMTTitleMapEntry{Décomposition équilibrée}{Décomposition unique \(p=r_3(p)+3c_3(p)\) en résidu équilibré et retenue}% ADP102-0307
  \HMTTitleMapEntry{Involution}{Action de \((\tau_+,\tau_-)\mapsto(-\tau_-,-\tau_+)\) sur les caractères \(p\) et \(a\)}% ADP102-0308
  \HMTTitleMapEntry{Définition et domaine}{Définition de \(L_p\) sur \(\mathbb R_{>0}\) et domaine de l'inverse \(\exp_p\)}% ADP102-0309
  \HMTTitleMapEntry{Addition déformée}{Loi de composition \(L_p(xy)=L_p(x)+L_p(y)+pL_p(x)L_p(y)\)}% ADP102-0310
  \HMTTitleMapEntry{Non-exhaustivité et portée exacte}{Domaine classificatoire de \(p\in\{2,1,0,-1,-2\}\) et extension affine générale}% ADP102-0311
  \HMTTitleMapEntry{Obstruction au lecteur réduit}{Défaut de naturalité du lecteur réduit \((p,a,\Lambda,C,m,r,\theta)\)}% ADP102-0312
  \HMTTitleMapEntry{Définition du lecteur complet}{Lecteur radial enrichi \(\mathcal P^{\mathrm{enr}}_{\mathrm{rad},N}\) : domaine, codomaine et conservation des préfixes}% ADP102-0313
  \HMTTitleMapEntry{Conservation comme traçabilité}{Conservation dynamique de la retenue, du bord et de la mémoire par le lecteur radial}% ADP102-0314
  \HMTTitleMapEntry{Subdivision quadruple}{Naturalité du lecteur sous le raffinement quadruple \(\operatorname{sd}_4\)}% ADP102-0315
  \HMTTitleMapEntry{Classification généalogique et non-fidélité de la publication}{Classification par la signature radiale pleine et défaut de fidélité de l'évaluation archimédienne}% ADP102-0316
  \HMTTitleMapEntry{Classification par 3 incréments visibles}{Classification cinématique au moyen des incréments \((\omega,\beta,\mu)\)}% ADP102-0317
  \HMTTitleMapEntry{3 fibres d'une même évaluation visible}{\(3\) classes généalogiques ayant une même image archimédienne}% ADP102-0318
  \HMTTitleMapEntry{Loi de vis}{Identité spectrale \(V^{-1}\mathscr S V=(2+\sqrt3)\mathrm i\mathscr S\) et équation de la suspension logarithmique}% ADP102-0319
  \HMTTitleMapEntry{Mode tritique commun}{Mode singulier commun des feuillets APP dans la carte canonique \(1,\ldots,9\)}% ADP102-0320
  \HMTTitleMapEntry{Ledger material}{Registre généalogique de la voie matérielle : extension survivante, signature entière et caractère adimensionnel}% ADP102-0321
  \HMTTitleMapEntry{Publication dimensionnelle ultérieure}{Réalisation de l'opérateur constitutif du vide dans les champs \((E_\sigma,B_\sigma)\)}% ADP102-0322
  \HMTTitleMapEntry{Courbe plane polaire}{Courbure et longueur d'arc d'une courbe polaire \(r=r(\theta)\)}% ADP102-0323
  \HMTTitleMapEntry{Spécialisations}{Évaluation de \(\kappa_{\mathrm F}\) dans les familles logarithmique, archimédienne, de Fermat, hyperbolique et lituus}% ADP102-0324
  \HMTTitleMapEntry{Suspension spatiale}{Immersion spatiale \(\Gamma(\theta)\) et critères différentiels de Frenet--Lancret}% ADP102-0325
  \HMTTitleMapEntry{Trichromie géométrique}{Coordination des régimes \(J^2=-I,0,+I\) avec la géométrie de Frenet}% ADP102-0326
  \HMTTitleMapEntry{Conclusion}{Théorème de fermeture des correspondances spirales généalogiques}% ADP102-0327
  \HMTTitleMapEntry{Produit semi-direct affine et composition des voies}{Monoïde affine \(\operatorname{End}(U_{\mathrm{APP}})\ltimes U_{\mathrm{APP}}\) et composition ordonnée des voies}% ADP102-0328
  \HMTTitleMapEntry{Tableau exhaustif de la carte locale à 2 orientations}{Énumération exhaustive de la carte locale \((\tau_+,\tau_-)\mapsto(p,a,r_3,c_3)\)}% ADP102-0329
  \HMTTitleMapEntry{5 régions de \texorpdfstring{\(\pi\)}{pi} et non-fidélité de l'évaluation}{Décomposition de la microfibre de \(\pi\) en 5 régions et multiplicité de l'évaluation archimédienne}% ADP102-0330
  \HMTTitleMapEntry{Algèbre puissance--logarithmique et 5 publications canoniques}{Algèbre de \(L_p\) et classification des 5 familles canoniques \(p\in\{2,1,0,-1,-2\}\)}% ADP102-0331
  \HMTTitleMapEntry{Dérivation des 5 courbes canoniques}{Évaluation de \(\exp_p(u_0+\beta m)\) dans les 5 familles canoniques}% ADP102-0332
  \HMTTitleMapEntry{Inversion des feuillets}{Conjugaison \(p\leftrightarrow-p\) des feuillets puissance--logarithmiques}% ADP102-0333
  \HMTTitleMapEntry{Secteur affine général}{Classification du secteur affine par la paire \((p,[\Lambda,C])\)}% ADP102-0334
  \HMTTitleMapEntry{Naturalité, équivariance et non-fidélité}{Naturalité sous troncature, équivariance nonadique et défaut de fidélité archimédienne}% ADP102-0335
  \HMTTitleMapEntry{Le lecteur de préfixes}{Lecteur de préfixes affines et conservation du registre enrichi}% ADP102-0336
  \HMTTitleMapEntry{Raffinements non triviaux}{Critère de compatibilité pour les raffinements généraux de voies}% ADP102-0337
  \HMTTitleMapEntry{Équivariance, non invariance}{Équivariance du cocycle radial sous l'holonomie \(\Gamma_9\)}% ADP102-0338
  \HMTTitleMapEntry{3 démonstrations de non-fidélité}{\(3\) obstructions explicites à la fidélité de l'évaluation archimédienne}% ADP102-0339
  \HMTTitleMapEntry{Même extrémité, ordre affine différent}{Coïncidence du produit final et dépendance de l'ordre des facteurs affines}% ADP102-0340
  \HMTTitleMapEntry{Même courbe, région généalogique différente}{Coïncidence de l'image archimédienne entre régions généalogiques distinctes}% ADP102-0341
  \HMTTitleMapEntry{Classification par niveau mathématique}{Stratification par domaine : histoire, lecteur, évaluation archimédienne et réalisation physique}% ADP102-0342
  \HMTTitleMapEntry{Classification par rotation, rayon et mémoire}{Classification par les incréments \((\omega,\beta,\mu)\)}% ADP102-0343
  \HMTTitleMapEntry{Réalisations ultérieures et contrôles de type}{Réalisations électromagnétique et matérielle avec séparation des domaines et codomaines}% ADP102-0344
  \HMTTitleMapEntry{Centre électronique et lecteur ultérieur}{Compatibilité du lecteur radial avec le point fixe électronique \((5,5)\) et l'atlas \(81/56/13/324\)}% ADP102-0345
  \HMTTitleMapEntry{Matrice fermée de contrôles négatifs}{Matrice des conditions de validité par niveau causal}% ADP102-0346
  \HMTTitleMapEntry{Première prolongation et frontière coinductive}{Prolongement coinductif initial de \(w_6\) jusqu'au bord de compatibilité d'ordre \(1\)}% ADP102-0347
  \HMTTitleMapEntry{Voie B : noyau de vacances \texorpdfstring{\(E_{21/22}\)}{E21/22}}{Racine simple \(\xi_9\) du noyau analytique d'ordre 9 et comparaison exacte \(\operatorname{Tr}_9(\xi_9^{-1})=\operatorname{Tr}_9(\alpha_{\mathrm{HMT}}^{-1})\)}% ADP102-0348
  % Purismo residual: rótulos activos localizados en el índice del sucesor.
  % Las sustituciones siguientes no alteran el nivel jerárquico ni el cuerpo.
  \HMTTitleMapEntry{Chronologie observable des deux curseurs}{Évolution chronologique de la paire de curseurs et détermination de l'état observable}% ADP102-0349
  \HMTTitleMapEntry{Émetteur fini à deux curseurs}{Émetteur fini associé à la paire de curseurs et génération d'états admissibles}% ADP102-0350
  \HMTTitleMapEntry{D'une cellule à deux coordonnées : naissance d'APP}{Construction d'APP par décomposition de 1 cellule en 2 coordonnées}% ADP102-0351
  \HMTTitleMapEntry{La branche de \(20\) blocs}{Prolongement de la section survivante au moyen de 20 blocs compatibles}% ADP102-0352
  \HMTTitleMapEntry{Du préfixe à un point réel}{Évaluation archimédienne du préfixe compatible et détermination du point réel}% ADP102-0353
  \HMTTitleMapEntry{Une émission reproductible de douze triades}{Émission reproductible de 12 triades et reconstruction de leur état enrichi}% ADP102-0354
  \HMTTitleMapEntry{Construction du caractère de clôture et des trajectoires généalogiques structurelles de \(\pi\)}{Construction du caractère de clôture et des trajectoires généalogiques de \(\pi\)}% ADP102-0355
  \HMTTitleMapEntry{La pente 1/5 imposée par les \(5\) régions}{Détermination de la pente 1/5 par le bilan des 5 régions orientées}% ADP102-0356
  \HMTTitleMapEntry{Fenêtre dodécaphasique de l'état signé et clôture réversible du cycle}{Restriction finie de l'état dodécaphasique signé et clôture réversible du cycle}% ADP102-0357
  \HMTTitleMapEntry{Fenêtre dodécaphasique finie}{Domaine fini de la clôture dodécaphasique}% ADP102-0358
  \HMTTitleMapEntry{Prolongement distant \texorpdfstring{\(\mathsf T_{\alpha,\infty}\)}{T alpha infini} de la voie entière}{Prolongement coinductif \(\mathsf T_{\alpha,\infty}\) de la clôture dodécaphasique entière}% ADP102-0359
  \HMTTitleMapEntry{Théorème des quatre canaux entiers}{Théorème de classification des 4 canaux entiers}% ADP102-0360
  \HMTTitleMapEntry{Premier encodeur intrinsèque et sa limite}{Codeur intrinsèque initial et existence de sa limite}% ADP102-0361
  \HMTTitleMapEntry{Hypothèses de typage et contrôles négatifs}{Conditions de typage et critères d'incompatibilité de la théorie des vacances nonadiques}% ADP102-0362
  \HMTTitleMapEntry{Voie chronologique de longueur, d'action et de spectre}{Composition chronologique de longueur, action et spectre dans la détermination de \(G\)}% ADP102-0363
  \HMTTitleMapEntry{Opérateur tordu et contrôle négatif déterministe}{Opérateur avec torsion et obstruction déterministe à la continuation multiplicative}% ADP102-0364
  \HMTTitleMapEntry{Stratification mathématique de l’objet de Majorana en quatre niveaux}{Stratification mathématique de l'objet Majorana en 4 niveaux}% ADP102-0365
  \HMTTitleMapEntry{Six caractères et quatre torsions de tresse}{Classification de 6 caractères par 4 classes de torsion de tresse}% ADP102-0366
  \HMTTitleMapEntry{obstructions et conditions spécifiques de non-dégénérescence}{Obstructions spécifiques et conditions de non-dégénérescence}% ADP102-0367
  \HMTTitleMapEntry{Quark haut et antiquark haut}{Secteurs quark \(u\) et antiquark \(\bar u\)}% ADP102-0368
  \HMTTitleMapEntry{Quark bas et antiquark bas}{Secteurs quark \(d\) et antiquark \(\bar d\)}% ADP102-0369
  \HMTTitleMapEntry{Contrôles de réfutation.}{Conditions de validité et d'incompatibilité de la réalisation équivariante des espèces}% ADP102-0370
  \HMTTitleMapEntry{Le carré local de \texorpdfstring{\(e\)}{e} et la rupture de mémoire}{Carré positionnel local de \(e\) et différenciation généalogique d'états ayant le même bloc visible}% ADP102-0371
  \HMTTitleMapEntry{La loi générale avant ses réalisations par espèce}{Antécédence causale de la Loi Générale des Masses par rapport aux réalisations par espèce}% ADP102-0372
  \HMTTitleMapEntry{Fondement unifié de la masse}{Factorisation extension–route et conservation du registre généalogique}% ADP102-0373
  \HMTTitleMapEntry{La bande de parcours}{Définition de la bande de route par l'orbite d'une section persistante}% ADP102-0374
  \HMTTitleMapEntry{Le canal d'incidence est une courbure locale, non un facteur auxiliaire}{Interprétation du canal d'incidence comme courbure locale du système de jauge}% ADP102-0375
  % Publicaciones terminales del capítulo 95: se conserva cada desarrollo
  % completo y se sustituye la enumeración narrativa por objeto, operación y resultado.
  \HMTTitleMapEntry{1re résolution : Riemann et la positivité comme complément construit}{Publication spectrale–polaire du continu par le complément positif de Guinand–Weil}% ADP102-0376
  \HMTTitleMapEntry{Existence, régularité et unicité globales pour Navier–Stokes par borne uniforme et compacité}{Publication solénoïdale du continu : borne coercive, compacité et régularité globale}% ADP102-0377
  \HMTTitleMapEntry{3e résolution : Yang–Mills et le saut du secteur de courbure}{Publication de jauge du continu : positivité par réflexion et saut du secteur de courbure}% ADP102-0378
  \HMTTitleMapEntry{4e résolution : Hodge et construction effective de la préimage algébrique}{Publication cohomologique du continu : relèvement compatible et construction effective de la préimage algébrique}% ADP102-0379
  \HMTTitleMapEntry{5e résolution : BSD et coïncidence des expressions analytique et arithmétique}{Publication déterminantielle du continu : comparaison Kato–Poitou–Tate et identité du terme principal de Birch–Swinnerton–Dyer}% ADP102-0380
  \HMTTitleMapEntry{L’identité de diffusion qui conserve la puissance croisée}{Identité de dispersion et conservation de la puissance croisée}% ADP102-0381
  \HMTTitleMapEntry{Persistance de l’objet terminal sous le passage à la limite}{Persistance de l'état terminal solénoïdal lors du passage à la limite}% ADP102-0382
  \HMTTitleMapEntry{Existence, régularité et unicité globale de la solution de Navier–Stokes obtenue par borne uniforme et compacité}{Théorème de clôture solénoïdale globale par borne uniforme et compacité}% ADP102-0383
  \HMTTitleMapEntry{L’opérateur universel d’extension}{Opérateur universel d'extension des classes cohomologiques compatibles}% ADP102-0384
  \HMTTitleMapEntry{Factorisation des entrelaceurs terminaux RH–NS–YM–Hodge–BSD à travers le continu généalogique commun}{Pertes d'information propres aux 5 publications classiques et conservation dans l'état enrichi}% ADP102-0385
  \HMTTitleMapEntry{Factorisation de RH–NS–YM–Hodge–BSD comme expressions terminales du continu généalogique commun}{Théorème de factorisation terminale de RH–NS–YM–Hodge–BSD à travers le continu généalogique commun}% ADP102-0386
  \HMTTitleMapEntry{Émergence corrélative de \texorpdfstring{\(5\)}{5} actions intrinsèques non exhaustives}{Décomposition expositive de l'action corrélative en 5 vues intrinsèques du même état}% ADP102-0387
  \HMTTitleMapEntry{Un même bord : flux intérieur et tension de jauge}{Factorisation commune de l'opérateur de frontière dans les réalisations solénoïdale et de jauge}% ADP102-0388
  \HMTTitleMapEntry{L'état commun de bord}{État enrichi de frontière et décomposition paritaire du transport}% ADP102-0389
  \HMTTitleMapEntry{L'opérateur minimal et son spectre}{Complexe de bord minimal et équivalence spectrale des laplaciens de degrés 0 et 1}% ADP102-0390
  \HMTTitleMapEntry{Lecture hydrodynamique}{Réalisation solénoïdale de l'opérateur de frontière : contraction du flot et conservation de l'énergie}% ADP102-0391
  \HMTTitleMapEntry{Lecture Yang--Mills}{Réalisation de jauge de l'opérateur de frontière : coercivité, quotient de jauge et saut spectral}% ADP102-0392
  % Promoción jerárquica de los certificados causales de \(\Gamma_9\) y de
  % la generación trítica de los canales enteros.
  \HMTTitleMapEntry{Filtration monodromique et calcul à profondeur prescrite}{Certification finie de Γ₉ : 396 niveaux, 2.376 trits, 43 tours complets, 387 paires \((K,K+9)\) par canal et fibre ambiante de 729 éléments}% ADP102-0393
  \HMTTitleMapEntry{Recensement fini des transitions et génération décimale}{Recensement exhaustif de 396 transitions, 2.376 trits, 43 tours et 387 paires de même phase}% ADP102-0394
  \HMTTitleMapEntry{Structure entière des lecteurs de constantes}{Classification arithmétique des canaux entiers et compatibilité avec la demi-couronne tritique}% ADP102-0395
}
\let\HMTNativeChapter\chapter
\let\HMTNativeSection\section
\let\HMTNativeSubsection\subsection
\let\HMTNativeSubsubsection\subsubsection
\let\HMTNativeParagraph\paragraph
\let\HMTNativeSubparagraph\subparagraph
\RenewDocumentCommand{\chapter}{s o m}{%
  \HMTResolvePublicTitle{#3}%
  \IfBooleanTF{#1}{\HMTNativeChapter*{\HMTResolvedTitle}}{%
    \IfNoValueTF{#2}{\HMTNativeChapter{\HMTResolvedTitle}}{\HMTNativeChapter[\HMTResolvedTitle]{\HMTResolvedTitle}}}}
\RenewDocumentCommand{\section}{s o m}{%
  \HMTResolvePublicTitle{#3}%
  \IfBooleanTF{#1}{\HMTNativeSection*{\HMTResolvedTitle}}{%
    \IfNoValueTF{#2}{\HMTNativeSection{\HMTResolvedTitle}}{\HMTNativeSection[\HMTResolvedTitle]{\HMTResolvedTitle}}}}
\RenewDocumentCommand{\subsection}{s o m}{%
  \HMTResolvePublicTitle{#3}%
  \IfBooleanTF{#1}{\HMTNativeSubsection*{\HMTResolvedTitle}}{%
    \IfNoValueTF{#2}{\HMTNativeSubsection{\HMTResolvedTitle}}{\HMTNativeSubsection[\HMTResolvedTitle]{\HMTResolvedTitle}}}}
\RenewDocumentCommand{\subsubsection}{s o m}{%
  \HMTResolvePublicTitle{#3}%
  \IfBooleanTF{#1}{\HMTNativeSubsubsection*{\HMTResolvedTitle}}{%
    \IfNoValueTF{#2}{\HMTNativeSubsubsection{\HMTResolvedTitle}}{\HMTNativeSubsubsection[\HMTResolvedTitle]{\HMTResolvedTitle}}}}
\RenewDocumentCommand{\paragraph}{s o m}{%
  \HMTResolvePublicTitle{#3}%
  \IfBooleanTF{#1}{\HMTNativeParagraph*{\HMTResolvedTitle}}{%
    \IfNoValueTF{#2}{\HMTNativeParagraph{\HMTResolvedTitle}}{\HMTNativeParagraph[\HMTResolvedTitle]{\HMTResolvedTitle}}}}
\RenewDocumentCommand{\subparagraph}{s o m}{%
  \HMTResolvePublicTitle{#3}%
  \IfBooleanTF{#1}{\HMTNativeSubparagraph*{\HMTResolvedTitle}}{%
    \IfNoValueTF{#2}{\HMTNativeSubparagraph{\HMTResolvedTitle}}{\HMTNativeSubparagraph[\HMTResolvedTitle]{\HMTResolvedTitle}}}}
% Las utilidades de degradación posteriores deben conservar el mapa de títulos.
\let\HMTBookChapter\chapter
\let\HMTBookSection\section
\let\HMTBookSubsection\subsection
\let\HMTBookSubsubsection\subsubsection
\let\HMTBookParagraph\paragraph
\let\HMTBookSubparagraph\subparagraph
\makeatother
